\documentclass[11pt,a4paper]{article}
\usepackage[T1]{fontenc}
\usepackage[utf8]{inputenc}
\usepackage[ngerman]{babel}
\usepackage{lmodern}
\usepackage{microtype}
\usepackage[left=2.5cm,right=2.5cm,top=2.3cm,bottom=2.3cm]{geometry}
\usepackage{setspace}
\setstretch{1.1}
\setlength{\parindent}{0pt}
\setlength{\parskip}{0.45em}
\usepackage[hidelinks]{hyperref}

\newcommand{\blltitle}{Titel der Besonderen Lernleistung}
\newcommand{\studentname}{Vorname Nachname}
\newcommand{\examdate}{[Datum des Kolloquiums]}

\begin{document}

{\Large\textbf{Thesenpapier zur Besonderen Lernleistung}}\\
\textbf{Titel:} \blltitle\\
\textbf{Name:} \studentname\\
\textbf{Kolloquium:} \examdate

\vspace{0.8cm}
\section*{Kernthesen}

\textbf{These 1:} [Klare, streitbare Kernaussage.]\\
\textit{Beleg aus der Arbeit:} [Kapitel/Ergebniszahl/Abbildung]

\textbf{These 2:} [Klare, streitbare Kernaussage.]\\
\textit{Beleg aus der Arbeit:} [Kapitel/Ergebniszahl/Abbildung]

\textbf{These 3:} [Klare, streitbare Kernaussage.]\\
\textit{Beleg aus der Arbeit:} [Kapitel/Ergebniszahl/Abbildung]

\textbf{These 4 (optional):} [Interdisziplinaere oder methodenkritische Zuspitzung.]\\
\textit{Beleg aus der Arbeit:} [Kapitel/Ergebniszahl/Abbildung]

\vspace{0.6cm}
\section*{Methodische Verteidigung in 3 Punkten}
\begin{enumerate}
\item Warum dieses Untersuchungsdesign? [Kurzbegruendung]
\item Warum diese Daten/Werkzeuge? [Kurzbegruendung]
\item Wichtigste Limitation und Umgang damit? [Kurzbegruendung]
\end{enumerate}

\section*{Diskussionsanker fuer das Pruefungsgespraech}
\begin{itemize}
\item [Moegliche kritische Rueckfrage 1 + argumentativer Antwortkern]
\item [Moegliche kritische Rueckfrage 2 + argumentativer Antwortkern]
\item [Moegliche kritische Rueckfrage 3 + argumentativer Antwortkern]
\end{itemize}

\end{document}
